% Generated draft: requires mathematical and source review.
\documentclass[11pt]{article}
\usepackage[margin=25mm]{geometry}
\usepackage{fontspec}
\setmainfont{Latin Modern Roman}
\usepackage{amsmath,amssymb,mathtools,unicode-math}
\setmathfont{Latin Modern Math}
\usepackage{enumitem,needspace,xurl,hyperref}
\setlength{\parindent}{0pt}
\setlength{\parskip}{5pt}
\setlength{\emergencystretch}{3em}
\newcommand{\N}{\mathbb{N}}
\newcommand{\Z}{\mathbb{Z}}
\newcommand{\Q}{\mathbb{Q}}
\newcommand{\R}{\mathbb{R}}
\newcommand{\C}{\mathbb{C}}
\newcommand{\F}{\mathbb{F}}
\newcommand{\A}{\mathbb{A}}
\newcommand{\PP}{\mathbb{P}}
\newcommand{\E}{\mathbb{E}}
\newcommand{\eps}{\varepsilon}
\newcommand{\dd}{\,\mathrm{d}}
\begin{document}
\section*{UnsolvedMath ID: 3100071}
\textbf{Source problem code:} catalogue entry
\subsection{Definitions and mathematical statement}
Write \(\F_2=\Z/2\Z\), and for \(n\ge1\) let \(d(n)=|\{a\in\N_{>0}:a\mid n\}|\). For a coefficient sequence \((b_n)_{n\ge0}\) over \(\F_2\), the proposed density is the natural density
\[
\lim_{N\to\infty}\frac1N|\{0\le n<N:b_n=1\}|=\frac1{32},
\]
including the assertion that this limit exists. Riasanovsky's question asks for this equality when the \(b_n\) are the coefficients of \(1/f\), with the positive-index coefficients of \(f\) equal to \(d(n)\bmod2\).

The constant coefficient is not specified, and this affects the meaning of the question. In fact,
\[
f_+(x)=\sum_{n\ge1}(d(n)\bmod2)x^n=\sum_{m\ge1}x^{m^2}
\]
has constant coefficient zero and leading term \(x\). It has no inverse in \(\F_2[[x]]\). It does have a Laurent-series inverse in \(\F_2((x))\); one possible precise question is the displayed density assertion for \(b_n=[x^n]f_+(x)^{-1}\), \(n\ge0\). Alternatively, assigning constant coefficient one gives
\[
f_0(x)=1+f_+(x),\qquad f_0^{-1}\in\F_2[[x]],
\]
and gives a different density assertion with \(b_n=[x^n]f_0^{-1}\). The source does not select between these conventions, so neither is silently adopted. The corrupted divisor-function symbol in the transcription is read here as the number-of-positive-divisors function described in words. The supplied research summary marks the specific \(1/32\) claim open and unverified; its informal theta-series and other analogies provide no precise normalization or resolution.
\subsection{Short English statement}
Does the reciprocal of the divisor-parity series over \(\F_2\) have coefficient \(1\) with natural density \(1/32\)? The constant coefficient and the intended formal-series ring must be clarified to obtain a unique statement.
\subsection{Sources}
\begin{itemize}
\item Local UnsolvedMath registry, permanent ID \(3100071\), exact source code \texttt{catalogue entry}; attributed to Riasanovsky.
\item Cooper, \emph{Combinatorial Problems I Like}, October 2020 snapshot, Classical Number Theory, source-order bullet 71; \url{https://people.math.sc.edu/cooper/combprob.html}. HTML extraction, recorded access 2026-07-29.
\item Supplied research summary, classified OPEN-TRIAGE, describes a divisor-parity/theta-series conjecture and says a resolution could not be verified. No precise primary citation or complete Riasanovsky reference is supplied.
\end{itemize}
\end{document}
